\documentclass[../main.tex]{subfiles}
\graphicspath{{\subfix{../img/}}{\subfix{../architectures images/}}}
\begin{document}

\chapter{Sprint 1 : Authentification et Gestion des Utilisateurs}

\section{Introduction}

Tout commence ici : rien d'autre dans TrustGym ne peut fonctionner sans un syst\`{e}me d'authentification robuste. Ce premier sprint pose ce socle --- double \acrshort{jwt} pour distinguer le personnel des membres, contr\^{o}le d'acc\`{e}s par r\^{o}les (\acrshort{rbac}) et gestion des modules \acrshort{saas}. C\^{o}t\'{e} mobile, un flux d'inscription en quatre \'{e}tapes, avec v\'{e}rification par code \acrshort{otp}, permet aux membres de rejoindre la plateforme sans friction.

\section{Sprint Backlog}

\begin{longtable}{|p{3cm}|p{2.5cm}|p{4cm}|p{5cm}|}
\hline
\textbf{Epic} & \textbf{Nature} & \textbf{User Story} & \textbf{Description} \\
\hline
\endfirsthead
\hline
\textbf{Epic} & \textbf{Nature} & \textbf{User Story} & \textbf{Description} \\
\hline
\endhead
\hline
\endfoot

\multirow{3}{3cm}{Authentification Admin} & Fonctionnalit\'{e} & En tant qu'administrateur, je veux me connecter via mon email et mot de passe & Impl\'{e}menter le endpoint \texttt{POST /api/auth/login} avec g\'{e}n\'{e}ration de token \acrshort{jwt}. \\
\cline{2-4}
& Fonctionnalit\'{e} & En tant qu'administrateur, je veux que mes routes soient prot\'{e}g\'{e}es & Mettre en place \texttt{JwtAuthGuard} et \texttt{JwtStrategy} pour la v\'{e}rification des tokens admin. \\
\cline{2-4}
& T\^{a}che & Configurer la strat\'{e}gie Passport.js & Int\'{e}grer Passport.js avec NestJS pour la gestion des gardes d'authentification. \\
\hline

\multirow{3}{3cm}{Authentification Membre} & Fonctionnalit\'{e} & En tant que membre, je veux me connecter depuis l'application mobile & Impl\'{e}menter \texttt{POST /api/auth/member/login} avec \texttt{MemberJwtGuard}. \\
\cline{2-4}
& Fonctionnalit\'{e} & En tant que membre, je veux v\'{e}rifier mon token au lancement & Cr\'{e}er \texttt{GET /api/auth/member/verify} pour la validation du token au d\'{e}marrage. \\
\cline{2-4}
& Fonctionnalit\'{e} & En tant que membre, je veux m'inscrire en plusieurs \'{e}tapes & Impl\'{e}menter le flux : v\'{e}rification email, confirmation d'identit\'{e}, code \acrshort{otp}, d\'{e}finition du mot de passe. \\
\hline

\multirow{3}{3cm}{Gestion des Utilisateurs} & Fonctionnalit\'{e} & En tant que SuperAdmin, je veux cr\'{e}er des utilisateurs avec des r\^{o}les & Impl\'{e}menter le \acrshort{crud} des utilisateurs avec attribution des r\^{o}les. \\
\cline{2-4}
& Fonctionnalit\'{e} & En tant qu'admin, je veux prot\'{e}ger les routes par r\^{o}le & Mettre en place \texttt{RolesGuard} pour la protection des routes. \\
\cline{2-4}
& Fonctionnalit\'{e} & En tant qu'admin, je veux contr\^{o}ler l'acc\`{e}s aux modules & Impl\'{e}menter \texttt{ModuleGuard} avec \texttt{SaasModuleEnum}. \\
\hline

\multirow{2}{3cm}{S\'{e}curit\'{e} Globale} & T\^{a}che & Configurer les pipes de validation & Mettre en place le \texttt{ValidationPipe} global et \texttt{TypeOrmExceptionFilter}. \\
\cline{2-4}
& T\^{a}che & Configurer Helmet et CORS & S\'{e}curiser les en-t\^{e}tes HTTP et configurer le partage cross-origin. \\
\hline

\caption{Sprint Backlog -- Sprint 1}
\label{tab:sprint1-backlog}
\end{longtable}

\section{Analyse}

\subsection{Diagramme de cas d'utilisation}

La figure~\ref{fig:sprint1-usecase} regroupe les cas d'utilisation de ce sprint : connexion admin, gestion des utilisateurs c\^{o}t\'{e} back-office, et inscription mobile membre.

\begin{figure}[H]
\centering
\resizebox{0.95\textwidth}{!}{%
\begin{tikzpicture}[
  uc/.style={draw, semithick, ellipse, fill=white, minimum width=2.8cm, minimum height=0.9cm, text width=2.6cm, align=center, font=\scriptsize},
  sys/.style={draw, semithick, rounded corners=5pt, inner sep=10pt},
  link/.style={semithick}
]

% System boundary
\node[sys, minimum width=6cm, minimum height=13cm] (system) at (4.5,0) {};
\node[font=\small\bfseries, anchor=south, yshift=2pt] at (system.north) {Sprint 1 -- Authentification};

% Use cases : single column at x=4.5, grouped by actor
\node[uc] (uc1)  at (4.5,  5.4) {Se connecter (web)};
\node[uc] (uc2)  at (4.5,  4.2) {V\'{e}rifier le token};

\node[uc] (uc3)  at (4.5,  3.0) {G\'{e}rer les utilisateurs};
\node[uc] (uc4)  at (4.5,  1.8) {Attribuer les r\^{o}les};
\node[uc] (uc5)  at (4.5,  0.6) {Configurer les modules SaaS};

\node[uc] (uc6)  at (4.5, -0.6) {S'inscrire (multi-\'{e}tapes)};
\node[uc] (uc8)  at (4.5, -1.8) {V\'{e}rifier le code OTP};
\node[uc] (uc7)  at (4.5, -3.0) {Se connecter (mobile)};
\node[uc] (uc9)  at (4.5, -4.2) {D\'{e}finir le mot de passe};
\node[uc] (uc10) at (4.5, -5.4) {R\'{e}initialiser le mot de passe};

% Actors (left), each centred on its own group
\stickfigure{(-2.5, 4.8)}{Admin / Caissier}
\draw[link] (-2.5, 4.8) -- (uc1.west);
\draw[link] (-2.5, 4.8) -- (uc2.west);

\stickfigure{(-2.5, 1.8)}{SuperAdmin}
\draw[link] (-2.5, 1.8) -- (uc3.west);
\draw[link] (-2.5, 1.8) -- (uc4.west);
\draw[link] (-2.5, 1.8) -- (uc5.west);

\stickfigure{(-2.5, -3.0)}{Membre}
\draw[link] (-2.5, -3.0) -- (uc6.west);
\draw[link] (-2.5, -3.0) -- (uc8.west);
\draw[link] (-2.5, -3.0) -- (uc7.west);
\draw[link] (-2.5, -3.0) -- (uc9.west);
\draw[link] (-2.5, -3.0) -- (uc10.west);

% Include relationship (routed on the right, outside the column)
\draw[dashed, -{Stealth[length=2mm]}] (uc6.east) to[bend left=55]
  node[right, font=\tiny] {$\ll$include$\gg$} (uc8.east);

\end{tikzpicture}
}%
\caption{Diagramme de cas d'utilisation -- Sprint 1}
\label{fig:sprint1-usecase}
\end{figure}

\subsection{Cas d'utilisation cibl\'{e}s}

Conform\'{e}ment au format d\'{e}fini en section~\ref{sec:uc_format}, ce sprint couvre : \textbf{Connexion administrateur} (Admin/Caissier), \textbf{Inscription membre} en quatre \'{e}tapes avec \texttt{$\ll$include$\gg$} sur la v\'{e}rification \acrshort{otp}, et \textbf{Gestion des utilisateurs} (SuperAdmin). Sp\'{e}cificit\'{e}s : l'inscription membre n\'{e}cessite que le membre ait \'{e}t\'{e} pr\'{e}-enregistr\'{e} par un administrateur ; \`{a} d\'{e}faut, l'\acrshort{api} renvoie une erreur invitant \`{a} contacter l'administration.

\subsection{Diagramme de s\'{e}quence : Authentification JWT}

La figure~\ref{fig:sprint1-seq} illustre le flux d'authentification pour un administrateur via le endpoint \texttt{POST /api/auth/login}.

\begin{figure}[H]
\centering
\resizebox{\textwidth}{!}{%
\begin{tikzpicture}

% Actors
\node[sdactor] (client) at (0,0) {Navigateur};
\node[sdactor] (frontend) at (4,0) {Next.js (NextAuth)};
\node[sdactor] (api) at (8.5,0) {NestJS API};
\node[sdactor] (guard) at (12.5,0) {LocalAuthGuard};
\node[sdactor] (db) at (16.5,0) {PostgreSQL};

% Lifelines
\draw[dashed] (client.south) -- ++(0,-10);
\draw[dashed] (frontend.south) -- ++(0,-10);
\draw[dashed] (api.south) -- ++(0,-10);
\draw[dashed] (guard.south) -- ++(0,-10);
\draw[dashed] (db.south) -- ++(0,-10);

% Messages
\draw[sdmessage] (0,-1.2) -- node[above, font=\scriptsize] {1: Saisir email + mot de passe} (4,-1.2);
\draw[sdmessage] (4,-2) -- node[above, font=\scriptsize] {2: POST /api/auth/login} (8.5,-2);
\draw[sdmessage] (8.5,-2.8) -- node[above, font=\scriptsize] {3: validate()} (12.5,-2.8);
\draw[sdmessage] (12.5,-3.6) -- node[above, font=\scriptsize] {4: SELECT user WHERE email} (16.5,-3.6);
\draw[sdreturn] (16.5,-4.4) -- node[above, font=\scriptsize] {5: User entity} (12.5,-4.4);
\draw[sdreturn] (12.5,-5.2) -- node[above, font=\scriptsize] {6: bcrypt.compare() OK} (8.5,-5.2);
\draw[sdmessage] (8.5,-6) -- node[above, font=\scriptsize] {7: jwtService.sign(payload)} (8.5,-6.8);
\node[font=\scriptsize, text width=3cm, align=center] at (8.5,-6.4) {};
\draw[sdreturn] (8.5,-7) -- node[above, font=\scriptsize] {8: \{token, user, role\}} (4,-7);
\draw[sdmessage] (4,-7.8) -- node[above, font=\scriptsize] {9: Stocker session NextAuth} (4,-8.6);
\draw[sdreturn] (4,-9) -- node[above, font=\scriptsize] {10: Redirection dashboard} (0,-9);

\end{tikzpicture}
}%
\caption{Diagramme de s\'{e}quence -- Authentification administrateur}
\label{fig:sprint1-seq}
\end{figure}

\section{R\'{e}alisation}

\subsection{Backend}

\subsubsection{Module d'Authentification}

Le \texttt{AuthenticationModule} g\`{e}re deux familles d'utilisateurs bien distinctes, avec des tokens \acrshort{jwt} \cite{jwt_rfc7519} s\'{e}par\'{e}s gr\^{a}ce \`{a} Passport.js \cite{passport} :

\begin{itemize}
  \item \textbf{Strat\'{e}gie Admin} (\texttt{JwtStrategy} + \texttt{JwtAuthGuard}) : utilis\'{e}e pour authentifier les administrateurs et caissiers via le endpoint \texttt{POST /api/auth/login}. Le token g\'{e}n\'{e}r\'{e} contient l'identifiant de l'utilisateur, son r\^{o}le, l'identifiant du partenaire et du centre.
  \item \textbf{Strat\'{e}gie Membre} (\texttt{JwtMemberStrategy} + \texttt{MemberJwtGuard}) : d\'{e}di\'{e}e aux membres de l'application mobile via \texttt{POST /api/auth/member/login}. Le token inclut l'identifiant du membre et ses informations de salle.
\end{itemize}

Le flux d'inscription mobile repose sur trois endpoints s\'{e}par\'{e}s :
\begin{enumerate}
  \item \texttt{POST /api/auth/member/check-email} : v\'{e}rifie l'existence du membre et envoie un code \acrshort{otp} \`{a} six chiffres par email.
  \item \texttt{POST /api/auth/member/verify-email} : valide le code \acrshort{otp} saisi par le membre.
  \item \texttt{POST /api/auth/member/register} : finalise l'inscription en hashant et stockant le mot de passe.
\end{enumerate}

\subsubsection{Module Utilisateur}

Le \texttt{UserModule} g\`{e}re le \acrshort{crud} des utilisateurs du back-office. L'entit\'{e} TypeORM \texttt{User} stocke les informations de connexion (\texttt{email}, \texttt{password} hash\'{e}), le r\^{o}le (\texttt{roleId}), les r\'{e}f\'{e}rences au partenaire (\texttt{partnerId}) et \`{a} la salle (\texttt{gymId}), ainsi que les champs de synchronisation mat\'{e}rielle (\texttt{cardId}, \texttt{pin}, \texttt{isSynced}).

\subsubsection{Contr\^{o}le d'Acc\`{e}s}

Deux niveaux de contr\^{o}le d'acc\`{e}s ont \'{e}t\'{e} mis en place :

\begin{itemize}
  \item \textbf{\acrshort{rbac}} \cite{sandhu1996rbac} via \texttt{RolesGuard} : un d\'{e}corateur personnalis\'{e} \texttt{@Roles()} permet de restreindre l'acc\`{e}s \`{a} certaines routes selon le r\^{o}le de l'utilisateur authentifi\'{e}. Trois r\^{o}les sont d\'{e}finis : SuperAdmin, Admin et Cashier.
  \item \textbf{Contr\^{o}le modulaire} via \texttt{ModuleGuard} : le \texttt{SaasModuleEnum} d\'{e}finit les fonctionnalit\'{e}s disponibles par centre (Boutique, Coachs, Finance, Acc\`{e}s, \acrshort{kpi}, \acrshort{ai}), permettant une personnalisation par partenaire.
\end{itemize}

\subsubsection{Impersonation SuperAdmin}

Le syst\`{e}me d'impersonation permet au SuperAdmin de se connecter en tant qu'un autre utilisateur tout en pr\'{e}servant la tra\c{c}abilit\'{e} compl\`{e}te. Le token \acrshort{jwt} g\'{e}n\'{e}r\'{e} lors d'une session d'impersonation contient les champs suppl\'{e}mentaires \texttt{isImpersonation}, \texttt{impersonatedBy} et \texttt{impersonatorEmail}, garantissant que toute action effectu\'{e}e est attribu\'{e}e au v\'{e}ritable op\'{e}rateur dans la piste d'audit.

\subsubsection{Journal d'Audit des Connexions}

L'entit\'{e} \texttt{LoginLog} trace chaque connexion et d\'{e}connexion du back-office : horodatage, identifiant utilisateur, adresse IP. Les administrateurs consultent cet historique via \texttt{/m/login-logs} --- utile pour d\'{e}tecter une tentative suspecte ou v\'{e}rifier qui \'{e}tait connect\'{e} \`{a} un moment donn\'{e}.

\subsubsection{S\'{e}curit\'{e} Globale}

Plusieurs couches de protection sont configur\'{e}es directement dans \texttt{main.ts} :

\begin{itemize}
  \item \texttt{ValidationPipe} global avec transformation automatique et validation des \acrshort{dto}.
  \item \texttt{TypeOrmExceptionFilter} pour la gestion unifi\'{e}e des erreurs de base de donn\'{e}es (contraintes d'unicit\'{e}, r\'{e}f\'{e}rences manquantes).
  \item Helmet pour la s\'{e}curisation des en-t\^{e}tes HTTP (XSS, clickjacking, MIME sniffing).
  \item Configuration CORS restrictive pour autoriser uniquement les origines autoris\'{e}es.
  \item Pr\'{e}fixe global \texttt{/api} pour tous les endpoints.
\end{itemize}

\subsection{Frontend Web}

\subsubsection{Fournisseur de Session}

L'authentification c\^{o}t\'{e} web repose sur NextAuth.js, int\'{e}gr\'{e} via un \texttt{SessionProvider} dans la pile de fournisseurs racine. Le \texttt{JWT\_SECRET} partag\'{e} entre l'\acrshort{api} NestJS et NextAuth garantit la compatibilit\'{e} des tokens. L'\texttt{AuthContext} expose les informations de l'utilisateur connect\'{e} et les m\'{e}thodes de connexion/d\'{e}connexion \`{a} l'ensemble de l'application.

\subsubsection{Page de Connexion}

La page de connexion pr\'{e}sente un formulaire valid\'{e} par Yup \cite{yup} via React Hook Form. Apr\`{e}s soumission, le token \acrshort{jwt} est stock\'{e} dans la session NextAuth et l'utilisateur est redirig\'{e} vers le tableau de bord. La page g\`{e}re les erreurs d'authentification avec des messages d'erreur localis\'{e}s gr\^{a}ce \`{a} \texttt{next-intl} \cite{nextintl}.

\begin{figure}[H]
  \centering
  \includegraphics[width=\textwidth, height=0.75\textheight, keepaspectratio]{capture/Page de connexion web.png}
  \caption{Page de connexion du tableau de bord web}
  \label{fig:sprint1-web-login}
\end{figure}

\subsubsection{Gestion des Utilisateurs}

L'interface de gestion des utilisateurs offre une liste paginable avec recherche, cr\'{e}ation avec attribution de r\^{o}le, modification et d\'{e}sactivation. Le \texttt{ModuleAccessContext} contr\^{o}le la visibilit\'{e} des \'{e}l\'{e}ments d'interface selon les modules activ\'{e}s pour le centre.

\begin{figure}[H]
  \centering
  \includegraphics[width=\textwidth, height=0.75\textheight, keepaspectratio]{capture/Gestion des utilisateurs superadmin.png}
  \caption{Interface de gestion des utilisateurs du back-office}
  \label{fig:sprint1-user-management}
\end{figure}

\subsection{Frontend Mobile}

\subsubsection{Flux d'Inscription Multi-\'{E}tapes}

L'inscription sur l'application mobile suit un parcours en quatre \'{e}tapes, chacune correspondant \`{a} un \'{e}cran d\'{e}di\'{e} dans le dossier \texttt{app/auth/} :

\begin{enumerate}
  \item \textbf{SignUp} : saisie de l'adresse email et v\'{e}rification de son existence dans le syst\`{e}me via l'\acrshort{api}.
  \item \textbf{IdentityConfirm} : affichage et confirmation des informations personnelles du membre (nom, pr\'{e}nom).
  \item \textbf{VerifyEmail} : saisie du code \acrshort{otp} \`{a} six chiffres envoy\'{e} par email. Un timer de 60 secondes g\`{e}re le renvoi du code.
  \item \textbf{SetPassword} : d\'{e}finition du mot de passe avec validation (longueur minimale, complexit\'{e}).
\end{enumerate}

\begin{figure}[H]
  \centering
  \includegraphics[width=\textwidth, height=0.75\textheight, keepaspectratio]{capture/flux-inscription.jpeg}
  \caption{Flux d'inscription multi-\'{e}tapes sur l'application mobile}
  \label{fig:sprint1-mobile-signup}
\end{figure}

\subsubsection{\'{E}cran de Connexion Mobile}

L'\'{e}cran \texttt{SignIn} permet la connexion par email et mot de passe. Le formulaire est valid\'{e} par Yup \cite{yup}. Apr\`{e}s authentification r\'{e}ussie, le token \acrshort{jwt} est persist\'{e} dans \texttt{AsyncStorage} via le service \texttt{auth-storage.ts}, et l'utilisateur est redirig\'{e} vers les onglets principaux de l'application.

\begin{figure}[H]
  \centering
  \includegraphics[width=\textwidth, height=0.75\textheight, keepaspectratio]{capture/ecran-connexion-mobile.png}
  \caption{\'{E}cran de connexion de l'application mobile}
  \label{fig:sprint1-mobile-signin}
\end{figure}

\subsubsection{Journal des Connexions}

L'interface du journal des connexions (\texttt{/m/login-logs}) affiche l'historique complet des \'{e}v\'{e}nements d'authentification du back-office, permettant aux administrateurs de surveiller l'activit\'{e} de connexion de la salle.

\begin{figure}[H]
  \centering
  \includegraphics[width=\textwidth, height=0.75\textheight, keepaspectratio]{capture/Journal des connexions (Login Logs).png}
  \caption{Journal d'audit des connexions et d\'{e}connexions}
  \label{fig:sprint1-login-logs}
\end{figure}

\section{Conclusion}

Ce sprint pose les gardes-fous sur lesquels tout le reste s'appuie. Deux tokens \acrshort{jwt} distincts, un \acrshort{rbac} \`{a} trois niveaux, un contr\^{o}le modulaire par partenaire --- et, c\^{o}t\'{e} mobile, un onboarding en quatre \'{e}tapes qui guide le membre sans friction. Le Sprint 2 peut maintenant s'appuyer sur des identit\'{e}s v\'{e}rifi\'{e}es pour construire la gestion des membres et des abonnements.

\end{document}
